% basic text processing/formatting
\usepackage{ucs}                % more unicode
\usepackage[utf8x]{inputenc}
\usepackage{fontspec}
\usepackage[english]{babel}

% maths stuff
\usepackage{amsmath}
\usepackage{amsfonts}
\usepackage{amssymb}
\usepackage{mathtools}
\usepackage{slashed}            % (Feynman-) slashed symbols
\usepackage{commath}            % for nicer differentials
\usepackage{bm}                 % bold math (only if mathserif is set)
\usepackage{dsfont}             % double stroked characters

\usepackage{tikz}               % drawing all manners of things
\usepackage{datetime}           % for custom date and time format
\usepackage{array}              % for matrices in math environment
\usepackage{float}              % place graphics with "H"
\usepackage{braket}             % Dirac notation
\usepackage{placeins}           % FloatBarrier
\usepackage{listings}           % source code listings
\usepackage{lmodern}            % to scale font to whatever I want
\usepackage{verbatim}           % print input as it is
\usepackage[xparse, skins]{tcolorbox}   % nice coloured boxes
% for textblock, remove showboxes for release
\usepackage[absolute,overlay,showboxes]{textpos}

%% devel only
% overlay a grid over the slides
% \usepackage[texcoord,grid,gridcolor=black!10,subgridcolor=black!5,gridunit=mm]{eso-pic}
\usepackage{lipsum}


%----------------------------------------------------------
% general info
% ISO date format
\newdateformat{isodate}{\THEYEAR-\twodigit{\THEMONTH}-\twodigit{\THEDAY}}

% commands so we can refer to these later
\newcommand{\theauthor}{Jan-Lukas Wynen}
\newcommand{\theinstitute}{Institute for applied \LaTeX}
\newcommand{\thetitle}{Fancy shmancy talk}
\newcommand{\thesubtitle}{Does this subtitle appear anywhere?}
\newcommand{\thedate}{\isodate\today}
\newcommand{\thelogo}{\includegraphics[width=10mm]{logo.pdf}}

\author{\theauthor}
\title{\thetitle}
\date{\thedate}
\institute{\theinstitute}
\logo{\thelogo}


%----------------------------------------------------------
% setup beamer
\usetheme[titleinfoot, toc, rainbow]{aiphi}


%----------------------------------------------------------
% setup fonts
\newfontfamily\dvfamily{DejaVu Sans Mono}

\makeatletter
\def\verbatim@font{\dvfamily}
\makeatother


%----------------------------------------------------------
% tcolorbox
% box for itemizes
\NewTColorBox{itembox}{m}{
  title=#1,
  noparskip,  % not using library parskip
  left=1mm, right=3mm, bottom=3mm, top=3mm,  % less for left to account for itemize
  sharp corners,
  boxrule=0.5mm,
  colframe=aiphiblue,
  colback=white
}

% box with a fancy title
\NewTColorBox{fancybox}{m}{
  enhanced,
  title=#1,
  attach boxed title to top center={
    yshift=-0.25mm-\tcboxedtitleheight/2,  % extra shift is 1/2 boxrule
    yshifttext=1ex-\tcboxedtitleheight/2
  },
  boxed title style={ % draw wedges around the title
    boxrule=0.5mm,
    frame code={ \path[tcb fill frame] ([xshift=-4mm]frame.west)
      -- (frame.north west) -- (frame.north east) -- ([xshift=4mm]frame.east)
      -- (frame.south east) -- (frame.south west) -- cycle; },
    interior code={ \path[tcb fill interior] ([xshift=-2.5mm]interior.west)
      -- (interior.north west) -- (interior.north east)
      -- ([xshift=2.5mm]interior.east) -- (interior.south east) -- (interior.south west)
      -- cycle;}
  },
  noparskip,
  left=3mm, right=3mm, bottom=3mm,
  sharp corners,
  boxrule=0.5mm,
  colframe=aiphiviolet,
  colback=white,
  colbacktitle=aiphiviolet!70!white,
  coltitle=white
}

% box showing a tag besides the title
% arguments: title, tag, tag background, tag foreground
\NewTColorBox{tagbox}{m m O{black!80!white} O{white}}{
  title=#1,
  after title={\hfill\colorbox{#3}{\color{#4}#2}},
  boxsep=0.5mm,
  noparskip,
  left=3mm, right=3mm, bottom=3mm, top=3mm,
  sharp corners,
  boxrule=0.5mm,
  colframe=black!80!white,
  colback=white,
  colbacktitle=white,
  coltitle=black
}

% box to cite stuff in
\NewTColorBox{icitebox}{m}{
  enhanced,
  hbox,  % shrink to horizontal width of content
  nobeforeafter,  % place in line with text
  tcbox raise base,  % adjust vertical position
  boxrule=0.4pt,
  top=0mm, bottom=0mm, right=1.5pt, left=0mm,  % most margins controlled through title box
  arc=0pt,
  boxsep=-0.35pt,  % negative to remove gap between title colorbox and rule
  colframe=aiphigreen!80!white,
  coltext=aiphigreen!90!white,
  colback=aiphigreen!10!white,
  fontupper=\scriptsize\ttfamily,
  fonttitle=\scriptsize\ttfamily,
  attach title to upper,  % show title at left instead of above
  title={\setlength{\fboxsep}{1.7pt}\colorbox{aiphigreen!90!white}{#1}},  % extra box to give background to title
}

% a color box inside a textblock to allow for arbitrary positioning
\newenvironment{aside}[2]{
  \begin{textblock*}{#1} (#2)
    \begin{tcolorbox}[
      enhanced,
      noparskip,  % not using library parskip
      left=1mm, right=1mm, bottom=1mm, top=1mm,
      sharp corners,
      boxrule=0.01mm,
      colback=aiphiyellow,
      coltext=black!80!white,
      watermark tikz={\draw[line width=2mm] circle (1cm) node{\fontfamily{ptm}\fontseries{b}\fontsize{20mm}{20mm}\selectfont!};},
      ]
    }
    {\end{tcolorbox}\end{textblock*}}

% ----------------------------------------------------------
% TikZ
\usetikzlibrary{tikzmark, calc, arrows, spy}


% ----------------------------------------------------------
% code listings

% Define a new language for C++ to have proper detection of built in datatypes.
% I hope I got all keywords.
\lstdefinelanguage{aiphicpp}{
  morekeywords={class, struct, template, typename, decltype, sizeof, static, if, else, for, while, do, goto, static_cast, dynamic_cast, reinterpret_cast, const_cast, constexpr, const, inline, and, or, xor, bitand, bitor, and_eq, or_eq, xor_eq, return, throw, try, catch, noexcept, namespace, using, typedef, enum, union, beak, continue, switch, case, break, continue, default, delete, new, operator, public, private, protected, explicit, extern, mutable, volatile, override, virtual, final},
  morekeywords=[2]{void, unsigned, int, long, float, double, bool, char, auto},
  morecomment=[l]{//},
  morecomment=[s]{/*}{*/},
  morestring=[b]",
  morestring=[b]'
}

% Very simple language to highlight interactive shell commands, not scripts.
\lstdefinelanguage{aiphishell}{
  morekeywords={},
  morecomment=[l]{\#},
  morecomment=[s]{<}{>},
  morestring=[b]",
  morestring=[b]'
}

% custom style for code listings
\lstdefinestyle{aiphilst}{
  language=aiphicpp,  % the default language
  basicstyle=\scriptsize\dvfamily,  % default font size
  breakatwhitespace=false,  % if true, lines are only broken at whitespaces
  breaklines=true,
  breakindent=2em,
  postbreak=\mbox{\textcolor{aiphired}{\rotatebox[origin=c]{180}{$\Lsh$}}},
  showstringspaces=false,  % show marker at spaces in strings
  captionpos=b, % place captions below listing
%
  frame=l,                                  % frame around code, show only to the left
  framerule=1.5pt,                          % frame width
  rulecolor=\color{aiphibordeaux!75!white}, % colour of frame
%
  numbers=left,  % line numbers, one of none, left, right
  numbersep=6pt, % distance of line numbers from code
  numberstyle=\tiny\color{black},  % text style of line numbers
%
  commentstyle=\itshape\color{aiphipetrol},
  keywordstyle={[1]\bfseries\color{aiphiblue!75!white}},
  keywordstyle={[2]\color{aiphigreen}},
  stringstyle=\color{aiphired!90!black},
}
% set this style as default
\lstset{style=aiphilst}

\lstdefinestyle{aiphishell}{
  language=aiphishell,  % the default language
  basicstyle=\scriptsize\dvfamily,  % default font size
  breakatwhitespace=false,  % if true, lines are only broken at whitespaces
  breaklines=true,
  breakindent=2em,
  postbreak=\mbox{\textcolor{aiphiblue}{\rotatebox[origin=c]{180}{$\Lsh$}}},
  showstringspaces=false,  % show marker at spaces in strings
  captionpos=b, % place captions below listing
%
  frame=l,  % frame around code, show only to the left
  framerule=1.5pt,  % frame width
  rulecolor=\color{aiphiblue!75!white}, % colour of frame
%
  numbers=none,  % no line numbers here
%
  commentstyle=\itshape\color{aiphipetrol},
  keywordstyle={[1]\bfseries\color{aiphiblue!75!white}},
  keywordstyle={[2]\color{aiphigreen}},
  stringstyle=\color{aiphired!90!black},
}

% show a prompt in shell session listings
\makeatletter

% store language for shells as command so we can use it for comparison below
\newcommand\langname@aiphishell{}
\def\langname@aiphishell{aiphishell}

% define shell prompt
\newcommand\prompt@aiphishell{{\color{aiphiblue}\$}\ }

% add this to the beginning of every line (empty at first)
\newcommand\addedToEveryPar@aiphishell{}
\lst@AddToHook{EveryPar}{\addedToEveryPar@aiphishell}

% show the prompt if language is aiphishell
\lst@AddToHook{PreInit}{%
  \ifx\lst@language\langname@aiphishell%
    \let\addedToEveryPar@aiphishell\prompt@aiphishell%
  \fi
}

\makeatother

% environment to help with shell session listings
\lstnewenvironment{aiphishell}[1][]
{\lstset{language=aiphishell, style=aiphishell}}
{\lstset{language=aiphicpp, style=aiphilst}}


% ----------------------------------------------------------
% handle backup / appendix slides
% place these around backup slides so they are not counted by the progress bar
\newcommand{\backupbegin}{
   \newcounter{framenumberappendix}
   \setcounter{framenumberappendix}{\value{framenumber}}
}
\newcommand{\backupend}{
   \addtocounter{framenumberappendix}{-\value{framenumber}}
   \addtocounter{framenumber}{\value{framenumberappendix}}
}


%----------------------------------------------------------
% abbreviations
\newcommand{\unit}[1]{\,\text{#1}}
\newcommand{\ev}{\,\text{eV}}
\newcommand{\kev}{\,\text{keV}}
\newcommand{\mev}{\,\text{MeV}}
\newcommand{\gev}{\,\text{GeV}}

\renewcommand{\i}{\mathrm{i}}
\renewcommand{\epsilon}{\varepsilon}

\newcommand{\Tr}{\text{Tr}}


% ----------------------------------------------------------
% miscellanious

% remember current item command so we can switch back to it
\let\olditem\item

% cite things inline using an icitebox
\newcommand{\icite}[2]{\begin{icitebox}{#1}#2\end{icitebox}}
\newcommand{\arxiv}[1]{\icite{arXiv}{\href{https://arxiv.org/abs/#1}{#1}}}
\newcommand{\doi}[1]{\icite{doi}{\href{https://doi.org/#1}{#1}}}

%%% Local Variables:
%%% coding: utf-8
%%% mode: latex
%%% TeX-engine: xetex
%%% TeX-master: "main"
%%% End:
